\documentclass[11pt]{article}
\usepackage{mathblatt}
\usepackage{adjustbox,varwidth}
% gesetzt von werkzeuge/setzer.py aus dem Lernweg (L3-A · L3-B · L3-C · L3-D · L3-T) und den Bankzeilen; Muster T6B.
\input{praeambel}
\newcommand{\kennung}{FWE}
\newcommand{\grau}[1]{{\color{mbgrau}#1}}

\begin{document}
\pfheftstil
\fancyfoot[C]{\parbox[t]{\textwidth}{\mbfhbox\par\vspace{1.5mm}{\footnotesize\color{mbgrau}{\scriptsize \kennung}\hfill\thepage}}}
\pfheftkopf{Streifen- und Kreisdiagramm}{Klasse 7}
% L3-A: Streifendiagramm zeichnen
\begin{pfaufg}{}{1.}{}
\grau{a)\ So kommt die Klasse 7b zur Schule: Rad $40\,\%$, Bus $25\,\%$, Auto $10\,\%$, der Rest zu Fuß. Stelle das in einem Streifen dar, der $10$ cm lang ist.\par\smallskip Rest: $100\,\% - 40\,\% - 25\,\% - 10\,\% = 25\,\%$ gehen zu Fuß.\par Längen ($1\,\% = 1$ mm): Rad $4$ cm, Bus $2{,}5$ cm, Auto $1$ cm, zu Fuß $2{,}5$ cm.\par Kontrolle: $4 + 2{,}5 + 1 + 2{,}5 = 10$ cm.\par Von links nacheinander abmessen, Trennstrich ziehen, jeden Abschnitt beschriften.\par\streifenvoll{Rad/40,Bus/25,Auto/10,zu Fuß/25}}
\par\medskip
b)\ Ein Streifen ist $10$ cm lang und steht für $100\,\%$. Wie lang wird der Abschnitt?\par a)~$30\,\%$ \quad b)~$45\,\%$ \quad c)~$8\,\%$ \quad d)~$62\,\%$\karo{3}
\fusshilfe{\mbox{1:~b) \mbox{a) $3$ cm} \mbox{b) $4{,}5$ cm} \mbox{c) $0{,}8$ cm} \mbox{d) $6{,}2$ cm}}}
\end{pfaufg}

% L3-A: Streifendiagramm zeichnen
\begin{pfaufg}{}{2.}{}
Beim Schulfest wurden Getränke verkauft: Wasser $30\,\%$, Apfelsaft $25\,\%$, Kakao $15\,\%$, der Rest Eistee. Trage alle Anteile in den Streifen ein und beschrifte sie.\par\streifenleer\karo{3}
\fusshilfe{\mbox{2:~Eistee $30\,\%$}}
\end{pfaufg}

% L3-A: Streifendiagramm zeichnen
\begin{pfaufg}{}{3.}{}
Die Klassen 7a und 7b wurden nach dem Schulweg gefragt. 7a: Rad $45\,\%$, Bus $30\,\%$, der Rest zu Fuß. 7b: Rad $35\,\%$, Bus $50\,\%$, der Rest zu Fuß. Stelle jede Klasse in einem Streifen dar. In welcher Klasse ist der Anteil „zu Fuß“ größer? Woran siehst du es am Streifen?\par 7a\par\streifenleer[0]\par 7b\par\streifenleer[0]\karo{3}
\fusshilfe{\mbox{3:~7a ($25\,\%$ gegen $15\,\%$)}}
\end{pfaufg}

% L3-B: Mittelpunktswinkel berechnen
\begin{pfaufg}{}{4.}{}
\gzwei{\adjustbox{max width=\linewidth}{\begin{varwidth}{50cm}\kreissektor{108}{$108^\circ$}\end{varwidth}}}{\grau{a)\ Berechne den Mittelpunktswinkel. a)~Ein Sektor steht für $30\,\%$. b)~$12$ von $500$ Befragten antworten „weiß nicht“.\par\smallskip a) Anteil: $30\,\% = 0{,}30$. Winkel: $0{,}30 \cdot 360^\circ = 108^\circ$.\par b) Anteil: $12 : 500 = 0{,}024$.\par Winkel: $0{,}024 \cdot 360^\circ = 8{,}64^\circ \approx 8{,}6^\circ$ (auf eine Stelle nach dem Komma runden).}}
\par\medskip
b)\ Wie groß ist der Mittelpunktswinkel?\par a)~$25\,\%$ \quad b)~$50\,\%$ \quad c)~$10\,\%$ \quad d)~$75\,\%$\karo{3}
\fusshilfe{\mbox{4:~b) \mbox{a) $90^\circ$} \mbox{b) $180^\circ$} \mbox{c) $36^\circ$} \mbox{d) $270^\circ$}}}
\end{pfaufg}

% L3-B: Mittelpunktswinkel berechnen
\begin{pfaufg}{}{5.}{}
Wie groß ist der Mittelpunktswinkel?\par a)~$35\,\%$ \quad b)~$12\,\%$ \quad c)~$62{,}5\,\%$ \quad d)~$4{,}5\,\%$\karo{3}
\fusshilfe{\mbox{5:~\mbox{a) $126^\circ$} \mbox{b) $43{,}2^\circ$} \mbox{c) $225^\circ$} \mbox{d) $16{,}2^\circ$}}}
\end{pfaufg}

% L3-B: Mittelpunktswinkel berechnen
\begin{pfaufg}{}{6.}{}
Berechne den Mittelpunktswinkel.\par a)~$9$ von $24$ Kindern haben ein Haustier. \quad b)~$14$ von $400$ Losen gewinnen.\karo{3}
\fusshilfe{\mbox{6:~\mbox{a) $135^\circ$} \mbox{b) $12{,}6^\circ$}}}
\end{pfaufg}

% L3-B: Mittelpunktswinkel berechnen
\begin{pfaufg}{}{7.}{}
Bei einer Umfrage unter $540$ Personen antworten $14$ mit „Sonstiges“. Im Kreisdiagramm soll jeder Sektor mindestens $10^\circ$ groß sein, sonst passt keine Beschriftung hinein. Reicht der Platz für „Sonstiges“?\karo{3}
\fusshilfe{\mbox{7:~Nein, nur $\approx 9{,}3^\circ$}}
\end{pfaufg}

% L3-C: Kreisdiagramm zeichnen
\begin{pfaufg}{}{8.}{}
\gzwei{\adjustbox{max width=\linewidth}{\begin{varwidth}{50cm}\kreisdiagramm[1.1]{Rad/144,Bus/123,zu Fuß/93}\end{varwidth}}}{\grau{a)\ So kommen die Kinder der 7c zur Schule: $14$ mit dem Rad, $12$ mit dem Bus, $9$ zu Fuß. Zeichne ein Kreisdiagramm.\par\smallskip Alle zusammen: $14 + 12 + 9 = 35$ Kinder.\par Winkel: Rad $14 : 35 \cdot 360^\circ = 144^\circ$; Bus $12 : 35 \cdot 360^\circ \approx 123{,}4^\circ \approx 123^\circ$; zu Fuß $9 : 35 \cdot 360^\circ \approx 92{,}6^\circ \approx 93^\circ$. Zum Zeichnen auf ganze Grad runden.\par Kontrolle: $144^\circ + 123^\circ + 93^\circ = 360^\circ$.\par Geodreieck mit dem Nullpunkt auf $M$, die $0^\circ$-Linie auf den Radius nach oben. $144^\circ$ abtragen, Linie ziehen. Geodreieck drehen: Nullpunkt bleibt auf $M$, $0^\circ$-Linie auf die eben gezogene Linie legen. $123^\circ$ abtragen. Der Rest ist zu Fuß.\par Jeden Sektor beschriften.}}
\par\medskip
\gzwei{\adjustbox{max width=\linewidth}{\begin{varwidth}{50cm}\kreisleer[1.3]\end{varwidth}}}{%
b)\ Zeichne ein Kreisdiagramm mit drei Sektoren: $180^\circ$, $120^\circ$ und $60^\circ$. Beginne am Radius.\karo{3}}
\fusshilfe{\mbox{8:~b) Kreis mit $180^\circ$, $120^\circ$, $60^\circ$}}
\end{pfaufg}

% L3-C: Kreisdiagramm zeichnen
\begin{pfaufg}{}{9.}{}
\gzwei{\adjustbox{max width=\linewidth}{\begin{varwidth}{50cm}\kreisleer[1.3]\end{varwidth}}}{%
Getränke beim Sportfest: Wasser $42\,\%$, Saft $33\,\%$, der Rest Tee. Berechne die Winkel und zeichne das Kreisdiagramm.\karo{3}}
\fusshilfe{\mbox{9:~Wasser $151^\circ$, Saft $119^\circ$, Tee $90^\circ$}}
\end{pfaufg}

% L3-C: Kreisdiagramm zeichnen
\begin{pfaufg}{}{10.}{}
\gzwei{\adjustbox{max width=\linewidth}{\begin{varwidth}{50cm}\kreisleer[1.3]\end{varwidth}}}{%
Die $36$ Kinder der Theater-AG haben verschiedene Aufgaben: $15$ spielen mit, $9$ bauen die Bühne, $6$ nähen Kostüme, die übrigen machen Licht und Ton. Stelle das in einem Kreisdiagramm dar.\karo{3}}
\fusshilfe{\mbox{10:~$150^\circ$, $90^\circ$, $60^\circ$, $60^\circ$}}
\end{pfaufg}

% L3-D: Kreisdiagramme lesen und zuordnen
\begin{pfaufg}{}{11.}{}
\gzwei{\adjustbox{max width=\linewidth}{\begin{varwidth}{50cm}\kreisdiagramm[1.0]{A/30,B/90,C/20,D/60}\end{varwidth}}}{\grau{a)\ $200$ Befragte kommen so zur Arbeit: Auto $90$, Bahn $60$, Rad $30$, zu Fuß $20$. Welcher Sektor gehört zu welchem Verkehrsmittel? Wie groß ist der Winkel für zu Fuß?\par\smallskip Angaben nach Größe ordnen: Auto $90$, Bahn $60$, Rad $30$, zu Fuß $20$.\par Sektoren nach Größe ordnen: B, D, A, C. Also B Auto, D Bahn, A Rad, C zu Fuß.\par Prüfen: Auto ist $90$ von $200$, knapp die Hälfte – B ist fast ein Halbkreis. Bahn ist $60$ von $200$, also $30\,\%$, etwas mehr als ein Viertel – D ist etwas größer als ein rechter Winkel.\par Winkel für zu Fuß: $20 : 200 = 0{,}1$; $0{,}1 \cdot 360^\circ = 36^\circ$.}}
\par\medskip
\gzwei{\adjustbox{max width=\linewidth}{\begin{varwidth}{50cm}\kreisdiagrammleer[0.9]{Essen/30,Freizeit/50,Sparen/20}\end{varwidth}}}{%
b)\ Mia teilt ihr Taschengeld ein: Freizeit $50\,\%$, Essen $30\,\%$, Sparen $20\,\%$. Schreibe an jeden Sektor, wofür er steht.\karo{3}}
\fusshilfe{\mbox{11:~b) Halbkreis}}
\end{pfaufg}

% L3-D: Kreisdiagramme lesen und zuordnen
\begin{pfaufg}{}{12.}{}
\gzwei{\adjustbox{max width=\linewidth}{\begin{varwidth}{50cm}\kreisdiagramm[0.9]{A/25,B/50,C/15,D/10}\end{varwidth}}}{%
Schätze für jeden Sektor: Ist er die Hälfte, ein Viertel oder weniger als ein Viertel? Wie viel Prozent sind C und D zusammen?\karo{3}}
\fusshilfe{\mbox{12:~A Viertel, B Hälfte, C und D weniger}}
\end{pfaufg}

% L3-D: Kreisdiagramme lesen und zuordnen
\begin{pfaufg}{}{13.}{}
\gzwei{\adjustbox{max width=\linewidth}{\begin{varwidth}{50cm}\kreisdiagramm[1.1]{Duschen/168,A/96,B/30,C/126,D/60}\end{varwidth}}}{%
Eine Familie verbraucht an einem Tag $480$ Liter Wasser: Duschen $168$ l, Toilette $126$ l, Wäsche $96$ l, Geschirr $60$ l, Kochen und Trinken $30$ l. Im Kreisdiagramm ist nur „Duschen“ beschriftet. a)~Welcher Buchstabe gehört zu welchem Teil? b)~Wie groß ist der Winkel für Kochen und Trinken?\karo{3}}
\fusshilfe{\mbox{13:~a) C Toilette, A Wäsche, D Geschirr, B Kochen und Trinken}}
\end{pfaufg}

% L3-T: Probetest
\begin{pfaufg}{}{14.}{}
In einer Umfrage zur Freizeit nennen $20\,\%$ Sport, $15\,\%$ Musik, $45\,\%$ Medien, der Rest Freunde treffen. Stelle das Ergebnis in einem Streifendiagramm dar.\par\streifenleer[0]\karo{3}
\fusshilfe{\mbox{14:~Freunde $20\,\%$}}
\end{pfaufg}

% L3-T: Probetest
\begin{pfaufg}{}{15.}{}
$650$ Jugendliche wurden befragt. $18\,\%$ von ihnen lesen nie Zeitung. Wie groß ist der Mittelpunktswinkel für „nie“ im Kreisdiagramm?\karo{3}
\fusshilfe{\mbox{15:~$64{,}8^\circ$}}
\end{pfaufg}

% L3-T: Probetest
\begin{pfaufg}{}{16.}{}
\gzwei{\adjustbox{max width=\linewidth}{\begin{varwidth}{50cm}\kreisleer[1.6]\end{varwidth}}}{%
Wohin ging die letzte Urlaubsreise? Meer $46{,}5\,\%$, Berge $27{,}5\,\%$, Stadt $26\,\%$. Zeichne ein Kreisdiagramm.\karo{3}}
\fusshilfe{\mbox{16:~Meer $167^\circ$, Berge $99^\circ$, Stadt $94^\circ$}}
\end{pfaufg}

% L3-T: Probetest
\begin{pfaufg}{}{17.}{}
\gzwei{\adjustbox{max width=\linewidth}{\begin{varwidth}{50cm}\kreisdiagramm[1.4]{Kuchen/450,A/150,B/375,C/25,D/250}\end{varwidth}}}{%
Ein Schulfest bringt $1\,250$~€ ein: Kuchen $450$~€, Getränke $375$~€, Spiele $250$~€, Tombola $150$~€, Bastelstand $25$~€. Im Kreisdiagramm ist nur „Kuchen“ beschriftet. a)~Welcher Buchstabe gehört zu welchem Teil? b)~Lisa sagt: „Kuchen und Getränke bringen zusammen mehr als die Hälfte.“ Hat sie recht?\karo{3}}
\fusshilfe{\mbox{17:~a) B Getränke, D Spiele, A Tombola, C Bastelstand}}
\end{pfaufg}

\par\vfill\noindent{\footnotesize Vorher: Säulen-, Balken- und Liniendiagramme\quad\textbullet\quad Weiter: Kenngrößen\par}
\end{document}
